
\noindent
\begin{tabular}{ll}
\textbf{Thesis} & LLM-based Wheel Human Interaction Robot for Receptionist and Human Assistant \\
\textbf{Student} & Mr. Kirawut Chalermkitpaisan, Mr. Krittin Sakharin, \\
                 & Ms. Natcha Rungruang, Ms. Natwasa Manomaiwiboon \\
\textbf{Student ID.} & 65011333, 65011356, 65011386, 65011402 \\
\textbf{Degree} & Bachelor of Engineering \\
\textbf{Program} & Robotics and AI Engineering \\
\textbf{Year} & 2025 \\
\textbf{Thesis Advisor} & Asst. Prof. Dr. Sarucha Yanyong \\
%\textbf{Co-Thesis Advisor} & Assoc. Prof. Dr. Co-Thesis Advisor \\
\end{tabular}

\vspace{1cm}

\begin{center}
\textbf{\Large ABSTRACT IN ENGLISH}
\end{center}

\vspace{0cm}

This thesis presents the development of Satu, a wheeled robotic assistant created within the Robotics and AI Engineering program at King Mongkut's Institute of Technology Ladkrabang. The primary objective was to build an accessible robot capable of assisting with academic tasks, such as answering department-related questions in Thai and guiding users to specific locations within Building E-12. To achieve this, the system integrates a Large Language Model (LLM) supported by a knowledge retrieval mechanism for contextual conversation, while autonomous navigation is handled via a ROS2-based framework and pre-constructed floor maps.

The unified system coordinates speech processing, visual perception, and navigation across multiple embedded computing devices, all housed within a physical design optimized for a welcoming user experience. Evaluation through iterative testing showed consistent improvements in conversation understanding and navigation accuracy, with the final version meeting all predefined performance criteria. These findings suggest that the integration of conversational AI and mobile robotics provides a reliable and functional solution for real-world academic environments.